\chapter{Escrita Criativa, Storytelling e Conteúdo Longo}
\label{ch:criatividade}

\ficcao Na abertura de um conto, o leitor vê uma máquina, um cesto e uma meia. Na página seguinte, a protagonista jura que a outra foi roubada. Se o leitor já a viu, o interesse está em saber como a protagonista se enganou. Se ainda não a viu, tem uma pista por descobrir. A mesma meia pode servir duas histórias.

\section{Escolher o que o leitor sabe}

Experimenta três arranjos para essa cena. Mostra a pista antes da personagem e deixa o leitor antecipar a descoberta. Esconde-a no enquadramento e revela-a no momento em que muda a interpretação. Ou apresenta dois relatos compatíveis com os indícios e faz a personagem escolher entre eles. Cada arranjo muda o trabalho do leitor. O capítulo~\ref{ch:audiovisual} mostra o que acontece quando mudas o enquadramento sem mudar o cenário.

Antes de pedir prosa, escreve quem quer o quê, o que impede a personagem e a alteração que a cena precisa de produzir. Se a protagonista tem dez minutos para encontrar a meia antes de fechar a loja, sabes o que cada gesto custa. Um cronómetro basta para tornar algumas decisões interessantes.

\section{Uma ficha de voz para experimentar}

Usa exemplos próprios ou autorizados e descreve a voz com escolhas que consegues reconhecer. Uma narradora pode usar humor seco, prestar atenção aos objetos e evitar explicar a piada. A ficha orienta o rascunho. No texto, verifica o efeito de uma frase em vez de contar adjetivos à procura de uma fórmula.

\begin{promptlabbox}{Definir uma voz para uma cena}
\textbf{Objetivo e contexto:} Descrever traços de textos que possas utilizar e aplicá-los a uma cena fictícia.

Escreve uma ficha com o ponto de vista, duas escolhas de ritmo e um limite para a informação revelada. Pede uma cena de três parágrafos, ou escreve-a tu. Experimenta uma versão em que o leitor sabe mais do que a protagonista.

\textbf{Resultado esperado e validação:} Identifica o ponto em que o leitor recebe a pista e o que a personagem sabe nesse momento. Lê ambas as versões e localiza o que mudou.

\textbf{Falhas possíveis:} Copiar a redação das amostras, mudar o ponto de vista sem marcar a transição ou confundir uma voz reconhecível com uma boa cena.
\end{promptlabbox}

\section{Do plano à cena}

Para uma história curta, anota a situação inicial, a escolha e a consequência. Para uma série, acrescenta o que cada episódio resolve e o que transporta para o seguinte. A meio de um rascunho, podes descobrir que a escolha mais interessante está noutra cena. Atualiza o plano para que sirva o texto que estás a construir.

\begin{lstlisting}
Propõe três cenas a partir desta premissa.
Para cada cena, indica o objetivo da personagem,
a escolha que enfrenta e o que muda depois.
Mantém o ponto de vista e as regras fornecidas.
\end{lstlisting}

\begin{promptlabbox}{Comparar escrita direta e escrita com plano}
\textbf{Objetivo e contexto:} Verificar como duas formas de trabalhar resolvem a mesma premissa.

Fixa o público, a extensão aproximada e o acontecimento final. Escreve ou pede uma versão direta e outra com plano seguido de rascunho. Se usares um modelo, usa sessões separadas e regista chamadas e tempo. Faz a leitura sem o rótulo da condição quando for possível.

\textbf{Resultado esperado e validação:} Compara ligação entre escolhas e consequências, continuidade de ponto de vista e efeito do final. Regista pelo menos uma passagem que justifique cada observação. Aceita um empate quando os critérios não distinguirem as versões.

\textbf{Falhas possíveis:} Dar entradas ou tempo diferentes sem o registar, medir apenas comprimento ou generalizar a partir de um conto.
\end{promptlabbox}

\section{Personagens e continuidade}

Uma ficha de personagem pode ter desejo, medo, informação conhecida e uma contradição que apareça em ação. A gerente quer fechar a loja a horas, mas recusa deixar uma reclamação por resolver. Numa cena de diálogo, dá a cada pessoa uma razão diferente para ficar. O que alguém deixa por dizer pode ajudar a explicar a escolha que faz a seguir.

Para um universo maior, conserva uma ficha de regras, uma cronologia e o estado no fim de cada capítulo. Dá ao modelo os excertos que interessam à nova cena. Um erro típico é deixar uma personagem usar uma chave que entregou no capítulo anterior. Assinala a última posse confirmada e a cena em que houve a entrega. Para uma série factual, usa a mesma ideia com fontes, datas, definições e remissões. Uma alteração a uma definição pede revisão dos episódios que ainda a usam.

\begin{promptlabbox}{Auditar a continuidade de uma cena}
\textbf{Objetivo e contexto:} Confrontar uma cena com regras escritas antes do rascunho.

Escreve uma regra sobre a chave, uma sobre o horário e uma sobre o que a protagonista sabe. Cria duas cenas curtas, uma compatível e outra com uma violação deliberada. Pede que cada objeção indique a regra e o trecho. Podes fazer a auditoria à mão e usar o resultado como referência.

\textbf{Resultado esperado e validação:} A violação é localizada e a cena compatível fica sem contradições atribuídas. Regista críticas de estilo à parte.

\textbf{Falhas possíveis:} Inventar uma regra, ignorar uma mudança já mostrada ou tomar um flashback assinalado por erro de cronologia.
\end{promptlabbox}

\section{Uma história melhor e uma coleção mais parecida}

Doshi e Hauser atribuíram 293 participantes a condições diferentes ao escrever histórias de oito frases. Uns escreveram sem assistência. Outros podiam pedir uma ideia ou até cinco. As histórias das condições com acesso a ideias tiveram avaliações mais altas de novidade. Na comparação por \emph{embeddings}, também ficaram mais próximas do perfil médio das outras histórias da mesma condição~\citep{doshi2024}. Uma avaliação do conto e uma medida da variedade da coleção respondem a perguntas diferentes. O desenho do estudo incide em histórias curtas com ajuda na fase de ideação.

Se preparas uma coleção, revê as premissas, os conflitos e os finais lado a lado. Três cenas podem funcionar bem em separado e gastar a mesma surpresa quando lidas juntas. Experimenta conceber uma ideia antes de consultar o modelo. Depois escolhe que contribuição vale a pena conservar e regista porquê.
